% SIAM (SISC) version of the BIE spectral-distance paper; main article.
% Supplement: supplement.tex. Build: latexmk main.tex supplement.tex
% (latexmkrc adds the SIAM class directory). Page target: 20 pages
% including references.
\documentclass[review,hidelinks,onefignum,onetabnum]{siamart251216}
% Shared preamble of the SIAM (SISC) version: packages, macros, theorem
% environments, title and headers. Loaded by main.tex and supplement.tex.

\usepackage{amsfonts}
\usepackage{amssymb}
\usepackage{mathtools}
\usepackage{mathrsfs}
\usepackage{stmaryrd}
\usepackage{bm}
\usepackage{booktabs}
\usepackage{graphicx}
\usepackage{xcolor}
\usepackage{tikz}
\usepackage{pgfplots}
\pgfplotsset{compat=1.18}
\usepgfplotslibrary{groupplots}
\usepackage{algorithmic}

\newcommand{\creflastconjunction}{, and~}

\newsiamremark{remark}{Remark}
\newsiamremark{assumption}{Assumption}
\crefname{assumption}{Assumption}{Assumptions}

% Notation (same symbols as draft/bie-enclosure/preamble.tex).
\renewcommand{\Im}{\mathrm{Im}}
\renewcommand{\Re}{\mathrm{Re}}
\newcommand{\ii}{\mathrm{i}}
\newcommand{\ee}{\mathrm{e}}
\newcommand{\dd}{\mathrm{d}}
\newcommand{\CC}{\mathbb{C}}
\newcommand{\RR}{\mathbb{R}}
\newcommand{\ZZ}{\mathbb{Z}}
\newcommand{\jump}[1]{\llbracket #1\rrbracket}
\newcommand{\QP}{\mathrm{QP}}
\newcommand{\rin}{\rho_{\min}}
\newcommand{\rax}{\rho_{\max}}

% Drafting note, printed in red; remove before submission.
\newcommand{\todo}[1]{\textcolor{red}{[TODO: #1]}}

\headers{Guided-mode enclosures by layer potentials}{Author}

\title{Reliable enclosures of guided-mode eigenvalues in periodic
  waveguides with a line defect by boundary integral equations\thanks{Submitted
  to the editors DATE.}}

\author{Author\thanks{Department of Mathematics, Example University, Example City
  (\email{author@example.com}).}}

\usepackage{amsopn}
\DeclareMathOperator{\diag}{diag}
% Cross-document references cannot be hyphenated; allow looser lines.
\setlength{\emergencystretch}{1.5em}

% xr-hyper appends the external file name to every imported label, including
% the two-field "@cref" labels of cleveref, which then prints it ("2.1main.pdf").
% Import "@cref" labels unchanged; all other labels keep the xr-hyper behavior.
\makeatletter
\long\def\XR@test#1#2#3#4\XR@{%
  \let\XR@tempa\@gobblefour
  \ifx#1\newlabel
    \let\XR@tempa\@firstoffour
  \else\ifx#1\XR@bibcite
    \let\XR@tempa\@secondoffour
  \else\ifx#1\@input
     \let\XR@tempa\@thirdoffour
  \else\ifx#1\new@label@record
     \let\XR@tempa\@fourthoffour
  \fi\fi\fi\fi
   \XR@tempa
    {%
     \in@{@cref}{#2}%
     \ifin@
       \expandafter\xdef\csname r@\XR@prefix#2\endcsname{\unexpanded{#3}}%
     \else
       \expandafter\protected@xdef\csname r@\XR@prefix#2\endcsname{\XR@addURL{#3}}%
     \fi
    }%
    {\expandafter\bibcite\expandafter{\XR@prefix#2}{#3}}%
    {\edef\XR@list{\XR@list\filename@area#2\relax}}%
    {%
      \edef\next{\noexpand\new@label@record{\XR@prefix#2}{\unexpanded{#3}{xr-url}{\XR@URL}}}%
      \next
    }
  \ifeof\@inputcheck\expandafter\XR@aux
  \else\expandafter\XR@read\fi}
\makeatother

\externaldocument[][nocite]{supplement}

\begin{document}
\maketitle

\begin{abstract}
Guided modes of a two-dimensional periodic waveguide with a line defect are
eigenvalues in a gap of the essential spectrum of a self-adjoint operator on an
unbounded strip. We compute candidates by boundary integral equations: lead
cells are coupled through Rayleigh amplitudes, decaying half-guide solutions
come from an ordered QZ decomposition of the cell scattering pencil, and the
center cell is reduced to a nonlinear eigenvalue problem for Cauchy data. From
a candidate we build a field on the whole strip that solves the Helmholtz
equation exactly in every homogeneous piece and jumps across rod boundaries,
artificial walls, and seams. We prove that the distance from the candidate to
the spectrum is bounded by the grouped $L^2$ norms of these jumps, weighted by
trace constants that are explicit in the geometry. All inputs, including the
infinitely many lead cells and a finite-dimensional test of the common gap, are
evaluated in ball arithmetic. For a star-shaped defect between two different
leads, one run of 14 minutes encloses a guided-mode eigenvalue in an interval
of width $2.7\times10^{-11}$; examples with an empty defect, different lead
periods, and modes near a gap edge give widths below $10^{-11}$.
\end{abstract}

\begin{keywords}
photonic crystal waveguide, guided mode, boundary integral equation,
eigenvalue enclosure, a posteriori error bound, ball arithmetic
\end{keywords}

\begin{MSCcodes}
65N25, 65N38, 65G20, 35P15, 78A50
\end{MSCcodes}

% Section 1. Budget 1.75 pages including title and abstract.
% Source: draft/bie-enclosure/s-introduction.tex (clean version).
\section{Introduction}\label{sec:intro}
A line defect in a two-dimensional photonic crystal breaks the periodicity
across the defect but keeps it along the defect. At frequencies in a band gap
of the surrounding crystal, the defect may support waves that propagate along
it and decay exponentially across it, the \emph{guided modes}
\cite{JoannopoulosJohnsonWinnMeade2008}. For a fixed Bloch wavenumber $\beta$
along the defect they are eigenvalues of a self-adjoint operator $A(\beta)$ on
an unbounded strip, located in a gap of its essential spectrum
\cite{Fliss2013}.

Two difficulties arise. The first is that the strip is unbounded across the
defect. Supercell methods truncate it with periodic conditions and incur a
model error that grows as the mode becomes less localized
\cite{Fliss2013,KlindworthSchmidtFliss2014}. Joly, Li, and Fliss constructed
exact boundary conditions for periodic half guides from cell problems
\cite{JolyLiFliss2006}, and Fliss used them to reduce the guided-mode problem to
a bounded neighborhood of the defect, at the price of a nonlinear eigenvalue
problem \cite{Fliss2013}; high-order finite element versions, Robin-to-Robin
variants, and dispersion-curve tracing followed
\cite{KlindworthSchmidtFliss2014,FlissKlindworthSchmidt2015,Klindworth2015}, and
related cell-DtN reductions are used in \cite{HuangLuLi2007,YuanLu2007}. The
second difficulty is reliability. For piecewise constant media, boundary
integral equations discretize only the interfaces and, with quasi-periodic Green
functions, reach high order with few unknowns
\cite{Linton1998,BarnettGreengard2010,HaoBarnettMartinssonYoung2014}. A nearly
singular matrix in a nonlinear eigenvalue problem, however, only shows that a
finite-dimensional equation is nearly solvable \cite{GuettelTisseur2017}; it
neither bounds the distance to the spectrum of the continuous operator nor
proves that a spectral point exists.

\emph{Related enclosure methods.} Bounding eigenvalues by the boundary defect of
trial fields that solve the equation exactly is the classical method of
particular solutions \cite{MolerPayne1968,BetckeTrefethen2005}. Descloux and
Tolley extended it to domain decomposition, with local exact solutions and an
eigenvalue error controlled by the mismatch of values and gradients between
subdomains \cite{DesclouxTolley1983} (see also \cite[Lem.~2.1]{Betcke2007});
this is the closest structure to ours. Barnett, Hassell, and Tacy bounded the
distance to the spectrum by boundary residuals for Neumann and Dirichlet
problems \cite{BarnettHassell2011,BarnettHassellTacy2018}. Dahne and Salvy
computed candidates by any means and then bounded the residual and the interior
norm with interval arithmetic and Taylor models \cite{DahneSalvy2020}; we
separate search and reliable evaluation in the same way. Temple, Kato--Temple,
and Lehmann--Goerisch bounds \cite{Temple1928,NakaoPlumWatanabe2019} require
trial functions in the domain of the operator or its form, which our
piecewise field is not. These methods do handle essential spectrum
\cite{ZimmermannMertins1995,DaviesPlum2004}, infinite waveguides
\cite{BehnkeMertinsPlumWieners2000}, and photonic-crystal line defects
\cite{BrownHoangPlumWood2009,HoangPlumWieners2009}; the finite element comparison
in \cref{sec:ex1} is of this type. All the particular-solution results above
concern bounded domains. Here the strip is unbounded, the discrete field is
generated cell by cell in the half guides, and besides the material interfaces
there are infinitely many artificial walls and seams.

\emph{Contributions.} We consider a strip with periodic media on both sides,
possibly different, and a local defect in a center cell, with piecewise
constant coefficients and one rod per cell. The main result is a
spectral-distance theorem for piecewise Helmholtz fields on this unbounded
strip (\cref{thm:distance}): for a trial field that solves the Helmholtz
equation exactly in every homogeneous piece,
$\operatorname{dist}(\widehat k^2,\sigma(A(\beta)))$ is bounded by the grouped
$L^2$ norms $\varepsilon_g$ of its jumps across rod boundaries and across the
infinitely many artificial walls and seams, weighted by trace constants $C_g$
that are explicit in the geometry. The resolvent mechanism behind such bounds
is classical; the theorem adapts it to a field generated cell by cell in both
half guides, with constants that can be evaluated. The second result is a
finite-dimensional test for a common gap of the two leads
(\cref{thm:complement,prop:panel}). A comparison with the Fourier complement
gives guaranteed lower bounds for the band above the gap from a small
plane-wave basis, and a Lipschitz bound in the Bloch wavenumber turns finitely
many inertia tests into bounds for all wavenumbers. Third, every input is
evaluated with rounding included: the jump integrals on curved interfaces, the
sums over infinitely many lead cells through the propagation matrix, the mass
lower bound, and the gap test. Together these give an interval that contains an
eigenvalue by the common gap alone, without a local isolation condition or a
mode count. For a star-shaped defect between two different leads, one run of
$14$ minutes gives a $\mu$-interval of width $2.74\times10^{-11}$; for modes
approaching a gap edge, where the lead contraction rate rises to $0.80$, the
width stays below $10^{-11}$.

Candidates are computed by boundary integral equations assembled from known
components: transparent boundary conditions for the half guides
\cite{JolyLiFliss2006,Fliss2013,FlissKlindworthSchmidt2015}, Rayleigh
expansions at artificial interfaces, and near-image splitting and proxy
representations of quasi-periodic fields
\cite{BarnettGreengard2010,BarnettGreengard2011,ChoBarnett2015}. None of these
is claimed as new, and the enclosure does not depend on the accuracy of the
candidate. We do not claim that the eigenvalue in the interval is unique, nor
its index or multiplicity, and we do not study a priori convergence: the
theorem applies to a given candidate and field, however they were obtained.

\Cref{sec:model} introduces the model and the gap test, \cref{sec:discretization}
the computation of candidates, and \cref{sec:bound} the trial field, the
spectral-distance theorem and the rounding-inclusive evaluation.
\Cref{sec:results} reports the examples, including a comparison with a finite
element enclosure, and \cref{sec:conclusion} concludes. The results stated here
are proved in the main text; the supplementary materials contain the error
bounds of the reliable evaluation and additional numerical data.

% Section 2. Budget 2.5 pages.
% Source: draft/bie-enclosure/s-modelproblem.tex (clean version).
\section{Model problem and the common gap}\label{sec:model}

\subsection{Periodic media with a line defect}
For $s\in\{-,+\}$ let $\rho_s$ be a coefficient with periods $L_s>0$ in $x$ and
$d>0$ in $y$; the two sides share the period $d$, while $L_-$ and $L_+$ may
differ. Each $\rho_s$ is piecewise constant: in a common background
$\rho_{\mathrm e}>0$, every periodic cell contains exactly one simply connected
rod with $C^2$ boundary and constant value $\rho_{s,\mathrm c}>0$. The
coefficient $\rho$ of the waveguide equals $\rho_-$ for $x<x_-$, $\rho_+$ for
$x>x_+$, and in the center cell $(x_-,x_+)\times I_d$, $x_\pm=\pm L_0/2$,
contains one further rod $\Omega_0$ or none. Here $I_d:=(-d/2,d/2)$, and the
cell windows are chosen so that every rod closure stays away from the cell
boundary. We write $\Upsilon$ for the union of all rod boundaries,
$0<\rin\le\rho\le\rax$, and $B:=\RR\times I_d$.

Fix $\beta\in[-\pi/d,\pi/d)$. A guided mode at wavenumber $k>0$ is a nonzero
solution of $-\Delta w=k^2\rho w$ in $\RR^2$ of the form
$w=v\,\ee^{\ii\beta y}$ with $v$ $d$-periodic in $y$ and $v|_B\in H^1(B)$
\cite{Fliss2013}; on $\Upsilon$ the value and the normal derivative of $w$ are
continuous. Let $H^1_\beta(B)$ be the functions in $H^1(B)$ whose traces on the
upper and lower edges $\Sigma^\uparrow,\Sigma^\downarrow$ satisfy
$u|_{\Sigma^\uparrow}=\ee^{\ii\beta d}u|_{\Sigma^\downarrow}$. On
$\mathcal H:=L^2(B,\rho\,\dd x\,\dd y)$ define, following
\cite[sect.~3]{Fliss2013},
\begin{equation}\label{eq:operator}
  A(\beta)u:=-\rho^{-1}\Delta u,\quad
  D(A(\beta)):=\bigl\{u\in H^1_\beta(B):\Delta u\in L^2(B),\
  \partial_y u|_{\Sigma^\uparrow}=\ee^{\ii\beta d}\partial_yu|_{\Sigma^\downarrow}\bigr\}.
\end{equation}
The condition $\Delta u\in L^2(B)$ excludes distributional sources on $\Upsilon$
and thus contains the transmission conditions. $A(\beta)$ is the nonnegative
self-adjoint operator of the closed form
$a[u,v]=\int_B\nabla u\cdot\overline{\nabla v}$, $D(a)=H^1_\beta(B)$, in
$\mathcal H$, and a guided mode is an eigenpair $A(\beta)u=k^2u$. Replacing
$\rho$ by $\rho_s$ in the whole strip gives the periodic operators $A_s(\beta)$
on $\mathcal H_s:=L^2(B,\rho_s\,\dd x\,\dd y)$. Since $\rho$ enters only the
mass term, Fourier analysis on the quasi-periodic cylinder gives
\begin{equation}\label{eq:domain-h2}
  D(A(\beta))=D(A_\pm(\beta))=H^2_\beta(B),
\end{equation}
with equivalent graph and $H^2$ norms. In particular, functions in
$D(A(\beta))$ have no jumps across $\Upsilon$; the trace estimate of
\cref{sec:bound} rests on this fact.

\subsection{Essential spectrum and the common gap}
For $\theta\in[-\pi/L_s,\pi/L_s)$ let $A_s(\beta,\theta)$ be the Bloch operator
on the cell $(0,L_s)\times I_d$ with phases $\ee^{\ii\theta L_s}$ in $x$ and
$\ee^{\ii\beta d}$ in $y$, and let $\lambda_{s,1}(\theta)\le\lambda_{s,2}(\theta)
\le\cdots$ be its eigenvalues. Here $\theta$ is the Bloch wavenumber of the lead
along $x$, not a second transverse phase: a band that crosses an energy at some
$\theta$ carries a wave into that lead.

\begin{proposition}[essential spectrum with two leads]\label{prop:ess}
For every real $\beta$,
\[
  \sigma_{\mathrm{ess}}(A(\beta))=\sigma(A_-(\beta))\cup\sigma(A_+(\beta)),
  \qquad
  \sigma(A_s(\beta))=\bigcup_{n\ge1}\bigcup_{\theta}\{\lambda_{s,n}(\theta)\}.
\]
\end{proposition}

For identical leads this is \cite[Prop.~3.1]{Fliss2013}; \cite[Rem.~2.1(2)]{Fliss2013}
notes the extension to different leads. A proof (a parametrix built from the two
periodic resolvents, and Weyl sequences supported far out in either lead) is
given in \cref{sm:proof-ess}. Since $\|\nabla u\|^2\ge\|\partial_y u\|^2\ge
b_\beta^2\|u\|^2_{L^2(B)}$ with $b_\beta:=\min_m|\beta+2\pi m/d|$, all three
operators are bounded below by $\lambda_0:=b_\beta^2/\rax$, which is positive
for $\beta\notin2\pi\ZZ/d$.

\begin{assumption}[common gap]\label{ass:gap}
For fixed $\beta$, the open interval $G_\beta=(g_-,g_+)$, $0<g_-<g_+$, does not
meet $\sigma(A_-(\beta))\cup\sigma(A_+(\beta))$.
\end{assumption}

\begin{proposition}[discrete spectrum in the gap]\label{prop:discrete}
Under \cref{ass:gap} the set $\sigma(A(\beta))\cap G_\beta$ consists of isolated
eigenvalues of finite multiplicity, which can accumulate only at $g_\pm$; the
eigenfunctions decay exponentially as $|x|\to\infty$.
\end{proposition}

The proof (\cref{sm:proof-discrete}) uses a conjugated resolvent of $A_\pm$
with a weight that vanishes on the support of the source. The spectral-distance
bound of \cref{sec:bound} holds for all of $\sigma(A(\beta))$; the gap is used
only to interpret the spectral point it produces.

\subsection{A finite-dimensional test for the common gap}\label{sec:gap-test}
To conclude that an interval $[\mu-R,\mu+R]$ contains an eigenvalue, we must
show for each lead that the $n_s$th band lies below the interval and the
$(n_s+1)$st band above it, for all $\theta$. Two ingredients reduce this to
finitely many matrix tests.

\begin{lemma}[Lipschitz bound for band roots]\label{lem:lipschitz}
If $\rho_s\ge\rho_{\min,s}>0$, then for all $n$, $\theta$, $\theta'$,
$|\lambda_{s,n}(\theta)^{1/2}-\lambda_{s,n}(\theta')^{1/2}|
\le|\theta-\theta'|/\rho_{\min,s}^{1/2}$.
\end{lemma}

\begin{proof}
After the gauge change $v=\ee^{-\ii\theta x}u$ the fiber form is
$a_\theta[v]=\|(\partial_x+\ii\theta)v\|^2+\|\partial_yv\|^2$ on a form domain
independent of $\theta$, and $|a_\theta[v]^{1/2}-a_{\theta'}[v]^{1/2}|\le
|\theta-\theta'|\,\|v\|_{L^2}\le|\theta-\theta'|\,\|v\|_{\mathcal H_s}/
\rho_{\min,s}^{1/2}$. Apply the min--max principle to the square root of the
Rayleigh quotient.
\end{proof}

The upper band needs a guaranteed \emph{lower} bound, which Ritz values do not
give. On the gauged cell take the orthonormal plane waves
$\psi_{p,q}=\ee^{2\pi\ii px/L_s}\ee^{\ii\beta_qy}/\sqrt{L_sd}$,
$\beta_q=\beta+2\pi q/d$, with symbols
$d_{p,q}(\theta)=(2\pi p/L_s+\theta)^2+\beta_q^2$. For a finite index set
$\mathcal I$ let $\Lambda_{\mathcal I}(\theta)=\inf_{(p,q)\notin\mathcal I}
d_{p,q}(\theta)$, $K_{\mathcal I}=\diag_{\mathcal I}d_{p,q}(\theta)$,
$M_{\mathcal I}=(\langle\rho_s\psi_j,\psi_i\rangle)_{i,j\in\mathcal I}$, and
$G_{\mathcal I}=(\langle\rho_s^2\psi_j,\psi_i\rangle)-M_{\mathcal I}^2$, the Gram
matrix of the components of $\rho_s\psi_j$ outside the span.

\begin{theorem}[Fourier complement comparison]\label{thm:complement}
Let $\lambda\ge0$ and $c_{\mathcal I}:=\Lambda_{\mathcal I}(\theta)-
\lambda\rho_{\max,s}>0$. With $T_{\mathcal I}:=K_{\mathcal I}-\lambda M_{\mathcal I}$
and $T^-_{\mathcal I}:=T_{\mathcal I}-\lambda^2G_{\mathcal I}/c_{\mathcal I}$, and
$n(\cdot)$ the number of negative eigenvalues,
\[
  n(T_{\mathcal I})\le\#\{j:\lambda_{s,j}(\theta)<\lambda\}\le n(T^-_{\mathcal I}).
\]
\end{theorem}

The left inequality is the Ritz bound; the right one follows from coercivity on
the complement after completing the square in the cross term
(\cref{sm:proof-complement}). Keeping the matrix $G_{\mathcal I}$, rather than a
scalar contrast bound, is what makes a lower bound possible with a small basis
($|\mathcal I|=169$ below). For a rod $\Omega_s$ with contrast
$\delta_s=\rho_{s,\mathrm c}-\rho_{\mathrm e}$ and
$B_{\mathcal I}=(\langle1_{\Omega_s}\psi_j,\psi_i\rangle)$ one has
$M_{\mathcal I}=\rho_{\mathrm e}I+\delta_sB_{\mathcal I}$ and
$G_{\mathcal I}=\delta_s^2(B_{\mathcal I}-B_{\mathcal I}^2)$; the entries of
$B_{\mathcal I}$ are Fourier coefficients of the rod, which reduce to periodic
boundary integrals.

\begin{proposition}[panel test]\label{prop:panel}
Cover $[-\pi/L_s,\pi/L_s]$ by panels with centers $\theta_i$ and radii $r_i$,
and let $0<k_L<k_U$, $\tilde r_i=r_i/\rho_{\min,s}^{1/2}<k_L$. If at every center
\begin{equation}\label{eq:panel-tests}
  n\bigl(T_{\mathcal I}((k_L-\tilde r_i)^2,\theta_i)\bigr)\ge n_s,\qquad
  n\bigl(T^-_{\mathcal I}((k_U+\tilde r_i)^2,\theta_i)\bigr)\le n_s,
\end{equation}
with $c_{\mathcal I}>0$ in the second test, then $\lambda_{s,n_s}(\theta)<k_L^2$
and $\lambda_{s,n_s+1}(\theta)\ge k_U^2$ for all $\theta$.
\end{proposition}

\begin{proof}
\Cref{thm:complement} gives $\lambda_{s,n_s}(\theta_i)<(k_L-\tilde r_i)^2$ and
$\lambda_{s,n_s+1}(\theta_i)\ge(k_U+\tilde r_i)^2$; by \cref{lem:lipschitz} the
band roots move by at most $\tilde r_i$ within a panel.
\end{proof}

The test also certifies the band index $n_s$. If the rod is symmetric under
$y\mapsto-y$, then $\lambda_{s,n}(\theta)=\lambda_{s,n}(-\theta)$ and the panels
need to cover only $[0,\pi/L_s]$. Writing $\overline a_s=k_L^2$ and
$\underline b_s=k_U^2$ for the certified values of lead $s$, the interval
$[\mu-R,\mu+R]$ lies in a common gap if
\begin{equation}\label{eq:gap-condition}
  \max(\overline a_-,\overline a_+)<\mu-R
  \quad\text{and}\quad
  \mu+R<\min(\underline b_-,\underline b_+).
\end{equation}
A rounding-inclusive implementation must also bound the quadrature error in
$B_{\mathcal I}$ (\cref{sm:band-tests}).

% Section 3. Budget 1.25 pages.
% Source: draft/bie-enclosure/s-reduction.tex (clean version).
\section{Reduction to the center cell}\label{sec:reduction}

Fix $\beta$. The center cell and the $j$th left and right periodic cells
($j\ge0$) are
\[
  \mathcal C_0=(x_-,x_+)\times I_d,\quad
  \mathcal C_{j,+}=\mathcal C_{0,+}+jL_+e_x,\quad
  \mathcal C_{j,-}=\mathcal C_{0,-}-jL_-e_x,
\]
with $\mathcal C_{0,+}=(x_+,x_++L_+)\times I_d$ and
$\mathcal C_{0,-}=(x_--L_-,x_-)\times I_d$; $B_\pm$ denotes the open half
strip formed by the cells on side $\pm$. The rods are $\Omega_0\subset\mathcal C_0$
(possibly empty) and $\Omega_{j,\pm}=\Omega_\pm\pm jL_\pm e_x$. The artificial
interfaces are $\Gamma_\pm=\{x_\pm\}\times I_d$ with normal $\nu_\pm=\pm e_x$,
pointing into $B_\pm$, and the walls $\Gamma_{j,\pm}=\{x_\pm\pm(j+1)L_\pm\}\times
I_d$ between consecutive cells. All normal derivatives on $\Gamma_\pm$ are taken
along $\nu_\pm$, from whichever side the function lives on.

For $O\in\{\mathcal C_0,B_-,B_+\}$ let $H^1_\beta(O)$ carry the quasi-periodic
condition in $y$, let $H^{1/2}_\beta(\Gamma_\pm)$ be its trace space and
$H^{-1/2}_\beta(\Gamma_\pm)$ the dual. For $u$ with $\Delta u\in L^2(O)$ and
quasi-periodic $\partial_yu$, the Cauchy data map
$\mathscr T_\pm u=(u|_{\Gamma_\pm},\partial_{\nu_\pm}u|_{\Gamma_\pm})$ is
defined weakly by Green's formula. The local solution spaces are
\begin{align*}
  \mathcal U_0(k)&=\{u\in H^1_\beta(\mathcal C_0):-\Delta u=k^2\rho u
  \text{ in }\mathcal C_0,\ \partial_yu\ \beta\text{-quasi-periodic}\},\\
  \mathcal V_\pm(k)&=\{v\in H^1_\beta(B_\pm):-\Delta v=k^2\rho_\pm v
  \text{ in }B_\pm,\ \partial_yv\ \beta\text{-quasi-periodic}\},
\end{align*}
with the equations in the distributional sense (so the rod transmission
conditions are included) and no condition on $\Gamma_\pm$. Functions in
$\mathcal V_\pm(k)$ have finite energy along the half strip. The
\emph{Cauchy data space} of the half guide is
$\mathcal G_\pm(k):=\mathscr T_\pm\mathcal V_\pm(k)$.

\begin{lemma}[unique continuation from Cauchy data]\label{lem:uc}
If $v\in\mathcal V_\pm(k)$ and $\mathscr T_\pm v=0$, then $v=0$.
\end{lemma}

\begin{proposition}[center reduction]\label{prop:reduction}
For fixed $\beta$ and $k>0$ let
$\mathcal U_0^{\mathcal G}(k)=\{u_0\in\mathcal U_0(k):\mathscr T_\pm u_0\in
\mathcal G_\pm(k)\}$. The restriction $u\mapsto u|_{\mathcal C_0}$ is a linear
bijection from $\ker(A(\beta)-k^2)$ onto $\mathcal U_0^{\mathcal G}(k)$.
\end{proposition}

The proofs are in \cref{sm:proof-reduction}: extending $v$ by zero across
$\Gamma_\pm$ through the homogeneous background strip gives an analytic
solution that vanishes on an open set; conversely, a center solution with
matching Cauchy data is glued to the unique half-guide solutions without a
distributional source on $\Gamma_\pm$.

\Cref{prop:reduction} holds for every $k>0$ and needs neither the gap nor
unique solvability of the half-guide Dirichlet problem. When the Dirichlet
problem is uniquely solvable, $\mathcal G_\pm(k)$ is the graph of the
Dirichlet-to-Neumann (DtN) operator and the reduction is the transparent
boundary condition of \cite[Thm.~4.5]{Fliss2013} (with the opposite normal).
At the countably many exceptional values where it is not
\cite[Thm.~4.1]{Fliss2013}, the DtN form fails while \cref{prop:reduction}
remains valid; Robin-to-Robin conditions avoid these values in another way
\cite{FlissKlindworthSchmidt2015}. In the gap, finite-energy half-guide
solutions decay cell by cell; Joly, Li, and Fliss represent them by a
propagation operator of spectral radius less than one \cite{JolyLiFliss2006},
whose discrete counterpart is the propagation matrix $\mathsf P_\pm$ of
\cref{sec:discretization}.

% Section 3. Budget 2.75 pages.
% Source: draft/bie-enclosure/s-discretization.tex (clean version).
\section{Discretization and candidate computation}\label{sec:discretization}

Candidates are computed on the center cell, as in the transparent boundary
conditions of Fliss \cite[Thm.~4.5]{Fliss2013}: a guided mode restricted to
$\mathcal C_0$ has port Cauchy data of finite-energy half-guide solutions. The
Dirichlet-to-Neumann form of this condition fails at countably many Dirichlet
eigenvalues in the gap \cite[Thm.~4.1]{Fliss2013}, which Robin-to-Robin
conditions avoid \cite{FlissKlindworthSchmidt2015}; we impose instead that the
port states lie in the decaying subspaces themselves. Lead cells are
discretized by boundary integral equations with port Rayleigh amplitudes, the
decaying half-guide solutions by an ordered QZ decomposition, and the center
cell by layer potentials plus background Rayleigh waves. The output is a
candidate $\widehat k$ with coefficients. We make no claim on its accuracy here:
\cref{sec:bound} bounds the spectral distance of whatever field these
coefficients define. Fix $\beta$ and a wavenumber $k>0$ with $k^2\in G_\beta$.

\subsection{Port Fourier basis and Rayleigh waves}
Let $\beta_m=\beta+2\pi m/d$, $\psi_m(y)=d^{-1/2}\ee^{\ii\beta_my}$, truncated to
$|m|\le M$, $K=2M+1$. With $\kappa_{\mathrm e}=k\rho_{\mathrm e}^{1/2}$ let
$\gamma_m=(\kappa_{\mathrm e}^2-\beta_m^2)^{1/2}$ with $\Im\gamma_m\ge0$, and
$\boldsymbol\gamma=\diag(\gamma_m)$. We assume that $k$ is not a Wood point,
that is, $\gamma_m(k)\ne0$ for all $m\in\ZZ$. Only the candidate search uses
this; the bound of \cref{sec:bound} does not.

In a vertical background strip, a $\beta$-quasi-periodic Helmholtz solution is
$u=\sum_m(a_mr_m^{x_p}+b_ml_m^{x_p})$ with
$r_m^{x_p}=\ee^{\ii\gamma_m(x-x_p)}\psi_m$ and
$l_m^{x_p}=\ee^{-\ii\gamma_m(x-x_p)}\psi_m$; $(a,b)$ are the right- and
left-going amplitudes anchored at $x_p$. Moving the anchor by $\ell$ multiplies
amplitudes by $\mathsf E(\ell)=\diag(\ee^{\ii\gamma_m\ell})$. For a port with
normal $\pm e_x$, amplitudes map to the Fourier coefficients of (value, normal
derivative) by $\mathsf C_\pm=\left[\begin{smallmatrix}I&I\\
\pm\ii\boldsymbol\gamma&\mp\ii\boldsymbol\gamma\end{smallmatrix}\right]$,
invertible away from Wood points.

\subsection{Layer potentials on a rod}\label{sec:layer}
Let $G^{(\kappa)}$ be the free-space Green function $\frac{\ii}{4}H^{(1)}_0(\kappa|
\mathbf z-\mathbf z'|)$ and $G^{(\kappa)}_\QP$ the $\ee^{\ii\beta d}$-quasi-periodic
one. We use the near-image split
\begin{equation}\label{eq:near-image}
  G^{(\kappa)}_\QP(\mathbf z,\mathbf z')=\sum_{n=-1}^{1}\ee^{\ii\beta nd}G^{(\kappa)}(\mathbf z,
  \mathbf z'+nd\,e_y)+G^{(\kappa)}_{\mathrm{far}}(\mathbf z-\mathbf z'),
\end{equation}
where $G^{(\kappa)}_{\mathrm{far}}$ is a smooth Helmholtz solution away from the
images $|n|\ge2$, since $(\Delta+\kappa^2)G^{(\kappa)}_\QP(\cdot,\mathbf 0)=
-\sum_n\ee^{\ii\beta nd}\delta_{nd\,e_y}$ \cite{Linton1998}. For a curve strictly inside a cell,
a quasi-periodic layer potential is thus three free-space layer potentials plus
a smooth correction near the closed cell, the structure used in
\cite{BarnettGreengard2010,BarnettGreengard2011,ChoBarnett2015}; the smooth
correction and the background are represented together by proxy point sources
outside the cell, with coefficients fitted to the quasi-periodicity on the top
and bottom seams. Since \cref{sec:bound} checks only the jumps of an explicit
field, the densities may be computed with any adequate approximation of
$G_\QP$; the field evaluated in \cref{sec:bound} uses the split
\eqref{eq:near-image} and fixed proxies, and all its mismatch enters its own
jumps.

For a rod $\Omega$ with $C^2$ boundary, outward normal $\nu$, and interior
wavenumber $\kappa_{\mathrm c}=k\rho_{\mathrm c}^{1/2}$, we represent the field in
the cell by
\begin{equation}\label{eq:cell-rep}
  u=\mathcal S^{(\kappa_{\mathrm e})}_\QP[\omega]+\mathcal D^{(\kappa_{\mathrm e})}_\QP[\tau]
  +u_{\mathrm b}\ \text{outside }\Omega,\qquad
  u=\mathcal S^{(\kappa_{\mathrm c})}[\omega]+\mathcal D^{(\kappa_{\mathrm c})}[\tau]
  \ \text{in }\Omega,
\end{equation}
with a finite background Rayleigh field $u_{\mathrm b}$. By the jump relations
\cite{ColtonKress1983,Kress1991}, the outside-minus-inside jumps of value and
normal derivative are, with $\eta=(\tau;-\omega)$,
\begin{equation}\label{eq:jump-operator}
  \mathscr J_k[\eta;u_{\mathrm b}]=
  \begin{bmatrix}
    I+D^{(\kappa_{\mathrm e})}_\QP-D^{(\kappa_{\mathrm c})}&-S^{(\kappa_{\mathrm e})}_\QP+S^{(\kappa_{\mathrm c})}\\
    T^{(\kappa_{\mathrm e})}_\QP-T^{(\kappa_{\mathrm c})}&I-D'^{(\kappa_{\mathrm e})}_\QP+D'^{(\kappa_{\mathrm c})}
  \end{bmatrix}\eta
  +\begin{bmatrix}u_{\mathrm b}\\ \partial_\nu u_{\mathrm b}\end{bmatrix}_{\partial\Omega}.
\end{equation}
The transmission conditions are $\mathscr J_k=0$, a second-kind system in which
the hypersingular parts cancel \cite{DominguezLyonTurc2016}. We discretize it by
Nystr\"om quadrature with Kress's logarithmic splitting on a uniform parameter
grid \cite{Kress1991,HaoBarnettMartinssonYoung2014}; $n$ is the number of density
unknowns. By the spectral representation of $G_\QP$, the potentials of a curve
lying left (right) of a line $x=x_p$ are purely right- (left-) going there;
their amplitudes are integrals of $\omega$ and $\tau$ against explicit
exponentials, and quadrature gives extraction matrices
$\mathsf F^{\rightarrow}_{x_p},\mathsf F^{\leftarrow}_{x_p}\in\CC^{K\times n}$.

\subsection{Lead cells: scattering matrix and decaying subspace}\label{sec:lead}
In the right cell $\mathcal C_{0,+}$ with ports $x_{\mathrm L}=x_+$,
$x_{\mathrm R}=x_++L_+$, the incident field
$u_{\mathrm b}=\sum_m(a_{\mathrm L,m}r_m^{x_{\mathrm L}}+b_{\mathrm R,m}l_m^{x_{\mathrm R}})$
determines $\eta$ through $\mathscr J_k[\eta;u_{\mathrm b}]=0$, and the outgoing
amplitudes are $b_{\mathrm L}=\mathsf E(L_+)b_{\mathrm R}+\mathsf F^{\leftarrow}_{x_{\mathrm L}}\eta$,
$a_{\mathrm R}=\mathsf E(L_+)a_{\mathrm L}+\mathsf F^{\rightarrow}_{x_{\mathrm R}}\eta$.
This defines the cell scattering matrix
$[b_{\mathrm L};a_{\mathrm R}]=[R_{\mathrm L},T_{\mathrm{RL}};T_{\mathrm{LR}},R_{\mathrm R}]
[a_{\mathrm L};b_{\mathrm R}]$, or equivalently, for the port states
$z_{\mathrm L}=(a_{\mathrm L};b_{\mathrm L})$ and $z_{\mathrm R}=(a_{\mathrm R};b_{\mathrm R})$,
\begin{equation}\label{eq:pencil}
  \mathsf E_{\mathrm L}z_{\mathrm L}=\mathsf E_{\mathrm R}z_{\mathrm R},\qquad
  \mathsf E_{\mathrm L}=\begin{bmatrix}-R_{\mathrm L}&I\\T_{\mathrm{LR}}&0\end{bmatrix},\quad
  \mathsf E_{\mathrm R}=\begin{bmatrix}0&T_{\mathrm{RL}}\\I&-R_{\mathrm R}\end{bmatrix}.
\end{equation}
The left cell, with its own rod, contrast and width $L_-$, is treated in the
same way. Floquet solutions satisfy $\mathsf E_{\mathrm L}z=\zeta\mathsf E_{\mathrm R}z$
with multiplier $\zeta=\ee^{\ii\theta L_+}$. In a gap the continuous
multipliers pair as $\zeta\leftrightarrow1/\bar\zeta$; the pairing
$\zeta\leftrightarrow1/\zeta$ of \cite{JolyLiFliss2006} corresponds to
replacing $\beta$ by $-\beta$ and agrees with the fixed-$\beta$ problem only
under $y\mapsto-y$ symmetry. Neither implies that the discrete pencil has
exactly $K$ multipliers inside the unit circle, so this is checked.

An ordered QZ decomposition $\mathsf Q\mathsf E_{\mathrm L}\mathsf Z=\mathsf U_{\mathrm L}$,
$\mathsf Q\mathsf E_{\mathrm R}\mathsf Z=\mathsf U_{\mathrm R}$ with the $K$
eigenvalues $|\zeta|<1$ first gives the leading columns $\mathsf Z_+$ and
\begin{equation}\label{eq:propagation}
  \mathsf P_+=\mathsf U_{\mathrm R,11}^{-1}\mathsf U_{\mathrm L,11},\qquad
  \mathsf E_{\mathrm L}\mathsf Z_+=\mathsf E_{\mathrm R}\mathsf Z_+\mathsf P_+ .
\end{equation}
Hence the port states $z_j=\mathsf Z_+\mathsf P_+^jc$ of consecutive cells are
compatible and decay geometrically: the discrete analogue of the
finite-energy solutions in the right half guide. On the left, the $K$ eigenvalues $1/\zeta$ inside the unit circle of the pencil
$(\mathsf E_{\mathrm R},\mathsf E_{\mathrm L})$ give $\mathsf Z_-$ and $\mathsf P_-$ with
$\mathsf E_{\mathrm R}\mathsf Z_-=\mathsf E_{\mathrm L}\mathsf Z_-\mathsf P_-$.
Using the whole Schur subspace keeps multiple and nonnormal multipliers; a
change of basis $\mathsf Z_\pm\mapsto\mathsf Z_\pm W$ only conjugates
$\mathsf P_\pm$. If the pencil is singular, the cluster dimension is not $K$,
or the cluster is not separated from the unit circle, the discretization at
that $k$ is rejected. The discrete Cauchy data spaces are
$\mathcal G^h_\pm(k)=\operatorname{range}\mathsf C_\pm\mathsf Z_\pm$.

\subsection{Center cell}\label{sec:center}
In $\mathcal C_0$ we use \eqref{eq:cell-rep} with $\Omega=\Omega_0$ (no layer
potentials if $\Omega_0=\varnothing$) and
\[
  u_{\mathrm b}=u_{\mathrm{hom}}=\sum_m\bigl(a^-_{0,m}r_m^{x_-}+b^+_{0,m}l_m^{x_+}\bigr),
\]
the part of the field entering through the two ports. It may be nonzero even
if $\Omega_0=\varnothing$, so quasi-periodic layer potentials alone do not
represent the center field. The ansatz follows the complementary-field
representation of \cite[App.~A]{BarnettGreengard2010} with the incident port
terms added. We use it only as an ansatz; neither completeness nor uniqueness of
the representation enters \cref{sec:bound}.

After truncation, the center unknown is $\xi=(\eta;a^-_0;b^+_0)\in\CC^{n+2K}$
($n=0$ for an empty center), and the discrete transmission equations read
$J(k)\xi=0$. The center field has the port amplitudes
\[
  \mathsf Y_-\xi=\begin{bmatrix}a^-_0\\ \mathsf F^{\leftarrow}_{x_-}\eta+\mathsf E(L_0)b^+_0\end{bmatrix},
  \qquad
  \mathsf Y_+\xi=\begin{bmatrix}\mathsf F^{\rightarrow}_{x_+}\eta+\mathsf E(L_0)a^-_0\\ b^+_0\end{bmatrix}.
\]

\subsection{Discrete eigenvalue problem and candidates}\label{sec:nep}
A candidate is a center solution whose port states lie in the decaying
subspaces: we ask for $\xi\ne0$ and $c_\pm\in\CC^K$ with
\begin{equation}\label{eq:nep}
  J(k)\xi=0,\qquad \mathsf Y_\pm\xi=\mathsf Z_\pm c_\pm .
\end{equation}
For root finding we need a scalar function that is analytic in $k$ and
independent of the basis of the stable subspaces; the determinant of the full
system \eqref{eq:nep} is neither. We therefore eliminate the center densities,
treat the center as a scattering cell with incident port amplitudes $\mathsf a\in\CC^{2K}$,
and let $\mathsf T_0(k)$, $\mathsf N_0(k)$ map $\mathsf a$ to the value and
normal-derivative coefficients of the total field at both ports. With the value
and normal-derivative blocks $\mathsf C_\pm\mathsf Z_\pm=[\mathsf D_\pm;\mathsf N_\pm]$
of the stable subspaces, the discrete DtN matrices are
$\Lambda_\pm=\mathsf N_\pm\mathsf D_\pm^{-1}$, and
\begin{equation}\label{eq:port-matrix}
  \mathsf A(k)=\mathsf N_0(k)-\diag\bigl(\Lambda_-(k),\Lambda_+(k)\bigr)\mathsf T_0(k)
  \in\CC^{2K\times2K}.
\end{equation}
On a complex neighborhood of a real interval free of Wood points, if both pencils
are regular with the stable cluster separated from the rest, and the matrices
used in the elimination are invertible, then $\det\mathsf A$ is analytic except
for poles where $\mathsf D_\pm$ or those matrices are singular. A singular
$\mathsf D_\pm$ is the discrete counterpart of a Dirichlet exceptional value
\cite[Thm.~4.1]{Fliss2013}, at which \eqref{eq:nep} remains valid. These conditions are monitored only at
interpolation nodes; they are empirical checks of the search, on which the final
interval does not depend.

\emph{Root search.} On each piece of the gap we interpolate $f(k)=\det\mathsf A(k)$
(computed by LU, as log-modulus and phase) at nested Chebyshev--Lobatto nodes of
degree $4,8,\dots,64$ until the trailing coefficients fall below $10^{-10}$ and two
off-grid checks agree to $10^{-9}$, bisecting otherwise; roots of the interpolant
come from the colleague matrix and are polished by Newton steps on the
interpolant \cite{Boyd2002,ZhaoBarnett2015}. A root is accepted if the scaled
singular-value ratio $s(\widehat k)=\sigma_{\min}/\sigma_{\max}$ of
$W_r\mathsf A(\widehat k)W_c$ is below $10^{-11}$; failures fall back to bisection
on $s(k)$. The candidate output is the right singular vector of the smallest
singular value, which determines $\mathsf a$, the center densities, and the lead
coordinates $c_\pm$. Together with $\widehat k$, $\mathsf P_\pm$, $\mathsf Z_\pm$
and the discretization parameters, these are the only inputs of
\cref{sec:bound}. A small $s(\widehat k)$ says that a finite system is nearly
singular, not how far $\widehat k^2$ is from $\sigma(A(\beta))$.

% Section 5. Budget 6 pages.
% Source: draft/bie-enclosure/s-error.tex (clean version).
\section{A spectral-distance bound for piecewise Helmholtz fields}\label{sec:bound}

Let $\mu:=\widehat k^2$. From the candidate we build a field $\widehat u$ on the
whole strip that solves the Helmholtz equation exactly in every homogeneous
piece but jumps across rod boundaries, artificial walls and the top/bottom
seam. We then prove $\operatorname{dist}(\mu,\sigma(A(\beta)))\le2\mu Q/(U-Q)$,
where $U=\|\widehat u\|_{\mathcal H}$ and $Q=\sum_gC_g\varepsilon_g$ combines
grouped jump norms $\varepsilon_g$ with trace constants $C_g$ that depend only on
the geometry, the coefficients and $\mu$. Nothing about how the candidate was
computed enters except through $\varepsilon_g$ and $U$.

\subsection{The trial field}\label{sec:field}
We separate \emph{definition parameters} (the density grid, the Rayleigh
truncation $M$, the candidate coefficients) from \emph{evaluation parameters}
(target quadrature, number of explicitly summed cells, evaluation precision).
Once the former are fixed, $\widehat u$ is a definite function; the latter only
affect how accurately $\varepsilon_g$ and $U$ are bounded. Both quantities are
homogeneous of degree one in the scaling of the candidate vector, so the bound
does not depend on its normalization.

\emph{Continuous densities.} On each rod boundary, parametrized by
$\boldsymbol r(t)$, $t\in[0,1)$, the $N_\eta$ density samples define the
trigonometric interpolant $\eta(t)=\sum_{m=-N_\eta/2}^{N_\eta/2-1}\eta_m
\ee^{2\pi\ii mt}$ through their discrete Fourier transform. All layer potentials
in $\widehat u$ are exact integrals of these trigonometric densities. The
Nystr\"om sums themselves are not used as a field: they carry point and dipole
singularities at the nodes and are not $H^2$ up to the interface.

\emph{Background representation.} In the background of each reference cell
$\mathcal C$ the field is
\begin{equation}\label{eq:mfs-field}
  \widehat u=\sum_{n=-1}^{1}\ee^{\ii\beta nd}\bigl(\mathcal S^{(\kappa_{\mathrm e})}[\omega]
  +\mathcal D^{(\kappa_{\mathrm e})}[\tau]\bigr)(\cdot-nd\,e_y)
  +\sum_{i=1}^{N_p}\alpha_iG^{(\kappa_{\mathrm e})}(\cdot,\mathbf p_i)+u_{\mathrm b},
\end{equation}
that is, the near-image split \eqref{eq:near-image} with the smooth correction
replaced by $N_p$ fixed proxy sources $\mathbf p_i$ outside the cell. Inside the
rod it is the free-space potential with $\kappa_{\mathrm c}$ as in
\eqref{eq:cell-rep}. We require that the image curves and the proxies have
positive distance from the closed cell and that all background terms use the
same $\kappa_{\mathrm e}$. The coefficients $\alpha$ are fitted to the whole field
after the densities are known: quasi-periodicity of value and $\partial_y$ at
collocation points on the seam lines, plus matching to outgoing Rayleigh waves
on two auxiliary vertical lines to exclude homogeneous solutions. The system
depends only on $\kappa_{\mathrm e}$, $\beta$ and the geometry, so one
factorization serves the center field and all lead states; for lead cells
$\alpha=\mathsf X_\pm v$ depends linearly on the state $v$. This is the
structure of \cite{BarnettGreengard2010}, except that the densities are not
re-solved together with $\alpha$. Saved floating-point coefficients are
treated as exact numbers; the accuracy of the fit, of the QZ decomposition, and
of the solver is not a hypothesis of the theorem, since every mismatch appears
in the jumps of this field.

\emph{Center and lead cells.} In $\mathcal C_0$, $\widehat u$ is
\eqref{eq:mfs-field} with $u_{\mathrm b}=u_{\mathrm{hom}}$ truncated to
$|m|\le M$ and $k=\widehat k$. In the right cell, incident amplitudes
$(a_{\mathrm L},b_{\mathrm R})$ give a density and hence a cell field
$u^{\mathrm{cell}}_+[a_{\mathrm L},b_{\mathrm R}]$, linear in the amplitudes.
Writing $\mathsf Z_+v=(a_+(v);b_+(v))$, set
$u_+[v]=u^{\mathrm{cell}}_+[a_+(v),b_+(\mathsf P_+v)]$ (left port state $v$,
right port state $\mathsf P_+v$), similarly $u_-[v]$, and
\begin{equation}\label{eq:lead-field}
  \widehat u(\mathbf z\pm jL_\pm e_x)=u_\pm[\mathsf P_\pm^jc_\pm](\mathbf z),
  \qquad\mathbf z\in\mathcal C_{0,\pm},\ j\ge0 .
\end{equation}
In exact arithmetic, \eqref{eq:propagation} makes neighboring cells agree on
the modes $|m|\le M$; the saved $\mathsf Z_\pm,\mathsf P_\pm$ satisfy this only
approximately, and the remaining mismatch, together with the modes $|m|>M$ and
the discrete residuals, appears as jumps on the shared walls.

The resulting \emph{trial field} satisfies:
(P1) $(-\Delta-\mu\rho)\widehat u=0$ exactly in every homogeneous piece;
(P2) the seam mismatch is treated as an interface jump;
(P3) for analytic rod boundaries, as in \cref{sec:results}, $\widehat u$ is $H^2$ up to
the boundary of each piece and has ordinary one-sided traces on walls and seams
(\cref{sm:proof-piecewise-h2});
(P4) $\widehat u\in\mathcal H$, which follows from the decay of $\mathsf P_\pm^j$ and
is checked in \cref{sec:computation}.
It is not in $H^1_\beta(B)$.

\subsection{Jumps and grouped defects}\label{sec:jumps}
The interfaces are the rod boundaries $\mathscr E_{\mathrm{mat}}$, the walls
$\mathscr E_{\mathrm{art}}=\{\Gamma_\pm,\Gamma_{j,\pm}\}$, and the seam segments
$\mathscr E_\Sigma$ of the cells, where the top and bottom edges of a cell are
glued (the field above $y=d/2$ is $\ee^{\ii\beta d}\widehat u(x,y-d)$). Each
shared curve is counted once. On $\Gamma$ choose side $1$ (the rod interior;
the cell nearer the center; $y<d/2$) and the normal $\nu_\Gamma$ from side $1$
to side $2$, and set
\begin{equation}\label{eq:jumps}
  \jump{\widehat u}_\Gamma=\widehat u^{(2)}-\widehat u^{(1)},\qquad
  \jump{\partial_\nu\widehat u}_\Gamma=\partial_{\nu_\Gamma}\widehat u^{(1)}
  -\partial_{\nu_\Gamma}\widehat u^{(2)} .
\end{equation}
On a seam this reads $\jump{\widehat u}_\Sigma=\ee^{\ii\beta d}\widehat u(x,-d/2)-\widehat u(x,d/2)$
and $\jump{\partial_\nu\widehat u}_\Sigma=\partial_y\widehat u(x,d/2)-\ee^{\ii\beta d}
\partial_y\widehat u(x,-d/2)$. On a rod boundary $\Gamma=\partial\Omega$ write the
outside field as $\mathcal S^{(\kappa_{\mathrm e})}[\omega]+\mathcal D^{(\kappa_{\mathrm e})}[\tau]+h$,
where $h$ (two images, proxies, $u_{\mathrm b}$) is smooth near $\Gamma$. The jump
relations give the \emph{actual} jumps
\begin{equation}\label{eq:actual-jumps}
\begin{aligned}
  \jump{\widehat u}_\Gamma&=\bigl(I+D^{(\kappa_{\mathrm e})}-D^{(\kappa_{\mathrm c})}\bigr)\tau
  +\bigl(S^{(\kappa_{\mathrm e})}-S^{(\kappa_{\mathrm c})}\bigr)\omega+h|_\Gamma,\\
  -\jump{\partial_\nu\widehat u}_\Gamma&=\bigl(T^{(\kappa_{\mathrm e})}-T^{(\kappa_{\mathrm c})}\bigr)\tau
  +\bigl(-I+D'^{(\kappa_{\mathrm e})}-D'^{(\kappa_{\mathrm c})}\bigr)\omega+\partial_\nu h|_\Gamma,
\end{aligned}
\end{equation}
with free-space operators on $\Gamma$ acting on the continuous densities. They
coincide with \eqref{eq:jump-operator} only if the proxy part of $h$ equals the
exact smooth correction; the difference is the fitting error, which stays in the
jumps. The reliable evaluation uses \eqref{eq:actual-jumps}.

We split the interfaces into finitely many groups $\mathscr E_g$; below there are
five: center rod, left rods, right rods, walls, seams. With a reference length
$\ell_g>0$ per group,
\begin{equation}\label{eq:defect}
  \varepsilon_g^2=\sum_{\Gamma\in\mathscr E_g}\Bigl(\|\jump{\partial_\nu\widehat u}_\Gamma\|^2_{L^2(\Gamma)}
  +\ell_g^{-2}\|\jump{\widehat u}_\Gamma\|^2_{L^2(\Gamma)}\Bigr),\qquad
  U=\|\widehat u\|_{\mathcal H}.
\end{equation}
Unlike a matrix residual, $\varepsilon_g$ includes the jumps between nodes, the
modes $|m|>M$, the seam mismatch, and all cells of both half guides.

\subsection{Grouped trace estimate}\label{sec:trace}
For each group take collars of half-width $\delta_g$ on both sides of every
interface: the normal map $(s,t)\mapsto\boldsymbol r_\Gamma(s)+t\nu_\Gamma(s)$,
$|t|<\delta_g$, is injective, the curvature satisfies $|\kappa_\Gamma|\le\kappa_g$
with $h_g=\delta_g\kappa_g<1$, and each point lies in at most $N_g$ collars of the
group. Collars are used only for functions in $D(A(\beta))=H^2_\beta(B)$, so they
may cross other interfaces, and collars of different groups may overlap.

\begin{proposition}[grouped trace estimate]\label{prop:trace}
Under these conditions let, for any $\eta_g>0$,
\[
  \tau_g=\frac{\delta_g}{6}+\frac{\delta_gh_g^2}{40(1-h_g^2)},\qquad
  D_{0,g}=\frac{N_g(1+\eta_g)}{2\delta_g},\qquad
  D_{1,g}=N_g(1+\eta_g^{-1})\tau_g .
\]
Then for $v\in D(A(\beta))$, with norms on $L^2(B)$ and
$\|D^2v\|^2=\|v_{xx}\|^2+2\|v_{xy}\|^2+\|v_{yy}\|^2$,
\begin{gather*}
  \sum_{\Gamma\in\mathscr E_g}\|v\|^2_{L^2(\Gamma)}\le D_{0,g}\|v\|^2+D_{1,g}\|\nabla v\|^2,\\
  \sum_{\Gamma\in\mathscr E_g}\|\partial_{\nu_\Gamma}v\|^2_{L^2(\Gamma)}\le D_{0,g}\|\nabla v\|^2
  +D_{1,g}\|D^2v\|^2 .
\end{gather*}
With $p_g(\lambda)=a_g+b_g\lambda+c_g\lambda^2$, $a_g=D_{0,g}/\rin$,
$b_g=D_{1,g}+\ell_g^2D_{0,g}$, $c_g=\ell_g^2D_{1,g}\rax$, and
\begin{equation}\label{eq:Cg}
  C_g^2=\sup_{\lambda\ge\lambda_0}\frac{p_g(\lambda)}{(\lambda+\mu)^2},
\end{equation}
the trace quantity
\[
  T_g(v)^2=\sum_{\Gamma\in\mathscr E_g}\bigl(\|v\|^2_{L^2(\Gamma)}
  +\ell_g^2\|\partial_{\nu_\Gamma}v\|^2_{L^2(\Gamma)}\bigr)
  \quad\text{satisfies}\quad T_g(v)\le C_g\|(A(\beta)+\mu)v\|_{\mathcal H}.
\]
\end{proposition}

The proof (\cref{sm:proof-trace}) applies a weighted one-dimensional trace
inequality on each normal segment, integrates with the Jacobian $1-t\kappa_\Gamma$,
and uses $\|D^2v\|_{L^2(B)}=\|\Delta v\|_{L^2(B)}$ on the quasi-periodic strip
(\cref{sm:global-h2}); the spectral theorem then turns
$\|v\|$, $\|\nabla v\|$, $\|\Delta v\|$ into the multiplier $p_g$. The supremum in
\eqref{eq:Cg} is attained at $\lambda_0$, at infinity (value $c_g$), or at the
stationary point $(2a_g-b_g\mu)/(2c_g\mu-b_g)$; if $2c_g\mu=b_g$ only the
endpoints matter. Write
\begin{equation}\label{eq:Q}
  Q=\sum_gC_g\varepsilon_g .
\end{equation}

\subsection{The spectral-distance theorem}\label{sec:theorem}
\begin{theorem}[spectral distance from interface defects]\label{thm:distance}
Assume the collar conditions of \cref{prop:trace}, $\mu>0$, and a nonzero field
$\widehat u$ with (P1)--(P4) and all $\varepsilon_g<\infty$. If $Q<U$, then
\[
  \operatorname{dist}\bigl(\mu,\sigma(A(\beta))\bigr)\le\frac{2\mu Q}{U-Q}.
\]
\end{theorem}

\begin{proof}
If $\mu\in\sigma(A(\beta))$ there is nothing to prove; let
$\delta=\operatorname{dist}(\mu,\sigma(A(\beta)))>0$ and, for the proof only,
$z=(A(\beta)-\mu)^{-1}\widehat u\in H^2_\beta(B)$. Write $u=\widehat u$.

\emph{Piecewise identity.} In each piece $\Delta u=-\mu\rho u$ and
$-\Delta z-\mu\rho z=\rho u$; since $\mu,\rho$ are real,
$\rho|u|^2=\overline z\Delta u-u\Delta\overline z$, and Green's second identity on
a bounded piece $D$ gives $\int_D\rho|u|^2=\int_{\partial D}(\overline z\,
\partial_{\nu_D}u-u\,\partial_{\nu_D}\overline z)\,\dd s$.

\emph{Shared interfaces.} Since $z\in H^2_\beta(B)$, its traces agree from both
sides of every interface, including the seam after quasi-periodic extension.
Adding the contributions of the two sides with outward normals $\pm\nu_\Gamma$
gives the integrand $\jump{\partial_\nu u}_\Gamma\overline z+\jump{u}_\Gamma
\partial_{\nu_\Gamma}\overline z$ on every $\Gamma\in\mathscr E$.

\emph{Truncation.} Let $\chi_X(x)$ be a smooth cutoff equal to $1$ for
$|x|\le X$ and $0$ for $|x|\ge2X$, with $\|\chi_X'\|_\infty=O(X^{-1})$ and
$\|\chi_X''\|_\infty=O(X^{-2})$. Replacing $z$ by $\chi_Xz$ involves finitely many
pieces and yields
\begin{multline*}
  \int_B\chi_X\rho|u|^2=\sum_{\Gamma}\int_\Gamma\chi_X\bigl(\jump{\partial_\nu u}_\Gamma\overline z
  +\jump{u}_\Gamma\partial_{\nu_\Gamma}\overline z\bigr)
  +\int_Bu\bigl(2\chi_X'\partial_x\overline z+\chi_X''\overline z\bigr)\\
  +\sum_\Gamma\int_\Gamma\jump{u}_\Gamma(\partial_{\nu_\Gamma}\chi_X)\overline z .
\end{multline*}
By weighted Cauchy--Schwarz and \cref{prop:trace},
\begin{equation}\label{eq:cs}
\begin{aligned}
  \sum_{\Gamma\in\mathscr E_g}\int_\Gamma\bigl(|\jump{\partial_\nu u}_\Gamma||z|
  +|\jump{u}_\Gamma||\partial_{\nu_\Gamma}z|\bigr)
  &\le\varepsilon_gT_g(z)\\
  &\le\varepsilon_gC_g\|(A(\beta)+\mu)z\|_{\mathcal H}<\infty .
\end{aligned}
\end{equation}
Hence the second term tends to zero ($u\in L^2$, $z\in H^2$), the third is at
most $\|\chi_X'\|_\infty\sum_g\ell_g\varepsilon_gT_g(z)\to0$, and dominated
convergence in the first gives
\[
  U^2=\sum_\Gamma\int_\Gamma\bigl(\jump{\partial_\nu u}_\Gamma\overline z+\jump{u}_\Gamma
  \partial_{\nu_\Gamma}\overline z\bigr)\,\dd s .
\]

\emph{Resolvent bound.} $\|z\|\le U/\delta$ and $(A+\mu)z=u+2\mu z$ give
$\|(A+\mu)z\|\le(1+2\mu/\delta)U$, so by \eqref{eq:cs}
$U^2\le Q(1+2\mu/\delta)U$. Since $U>Q$, $\delta(U-Q)\le2\mu Q$.
\end{proof}

The bound is independent of the scaling of $\widehat u$ and tends to zero
linearly as $\varepsilon_g\to0$. It holds for the whole spectrum, with no gap
hypothesis. A field with a volume residual $r=(-\Delta-\mu\rho)\widehat u$ adds
$\|\rho^{-1}r\|_{\mathcal H}$ to the numerator; ours has none.

\subsection{Computing the inputs}\label{sec:computation}
\emph{Half-guide sums.} By \eqref{eq:lead-field}, the jumps on
$\partial\Omega_{j,\pm}$, on the wall $\Gamma_{j,\pm}$ shared with the next cell,
and on the seam of $\mathcal C_{j,\pm}$ depend only on the state
$\mathsf P_\pm^jc_\pm$. Let $b_{g,\pm}(v)=\|\mathcal L_{g,\pm}v\|^2$ be the
group-$g$ part of these jump integrals for the state $v$, with $\mathcal L_{g,\pm}$
linear, and $b_{g,0}$ the group-$g$ part for the center cell (its rod, $\Gamma_\pm$
against $u_\pm[c_\pm]$, its seam). Then
\begin{equation}\label{eq:lead-sum}
  \varepsilon_g^2=b_{g,0}+\sum_\pm\sum_{j\ge0}b_{g,\pm}(\mathsf P_\pm^jc_\pm).
\end{equation}
The first $N_{\mathrm{cell}}$ terms (chosen per group and side) are evaluated
by iterating $v\leftarrow\mathsf P_\pm v$. If $b_{g,\pm}(v)\le\Theta_{g,\pm}\|v\|^2$ and
$\|\mathsf P_\pm^p\|_2\le r_\pm<1$ for some $p\ge1$, the tail satisfies
\begin{equation}\label{eq:tail}
  \sum_{j\ge N_{\mathrm{cell}}}b_{g,\pm}(\mathsf P_\pm^jc_\pm)
  \le\frac{\Theta_{g,\pm}\sum_{l<p}\|\mathsf P_\pm^l\|_2^2}{1-r_\pm^2}
  \bigl\|\mathsf P_\pm^{N_{\mathrm{cell}}}c_\pm\bigr\|_2^2 .
\end{equation}
The full nonnormal $\mathsf P_\pm$ is kept; $p>1$ is needed when
$\|\mathsf P_\pm\|_2\ge1$ although the spectral radius is below one. We bound
$\Theta_{g,\pm}$ by analytic envelopes of the kernels (\cref{sm:reliable-eps})
rather than by evaluating $K$ columns, and compensate the crude constant by a
slightly larger $N_{\mathrm{cell}}$.

\emph{Mass lower bound.} The same structure gives $U^2$ as a convergent sum, which
verifies (P4); the theorem needs only $U\ge(\int_D\rho|\widehat u|^2)^{1/2}$ for
any $D\subset B$. For a disk $D_r(\mathbf c)$ inside a piece with constant
$\rho_b$ and $\kappa_b=(\mu\rho_b)^{1/2}$, the circular mean of $\widehat u$ is
$J_0(\kappa_bs)\widehat u(\mathbf c)$, and Cauchy--Schwarz on circles gives
\begin{equation}\label{eq:disk}
  \int_{D_r(\mathbf c)}\rho|\widehat u|^2\ge\rho_b\pi r^2
  \bigl[J_0(\kappa_br)^2+J_1(\kappa_br)^2\bigr]|\widehat u(\mathbf c)|^2 ,
\end{equation}
so disjoint disks need only point values of the field.

\emph{Reliable evaluation.} Split each interface into parameter panels $I$ and let
$f(t)=|\boldsymbol r'(t)|^{1/2}(\jump{\partial_\nu\widehat u},
\ell_g^{-1}\jump{\widehat u})(\boldsymbol r(t))$. If a polynomial vector $P_I$
satisfies $\sup_I|f-P_I|\le e_I$, then
\begin{equation}\label{eq:panel}
  \|f\|_{L^2(I)}\le V_I^{1/2}+|I|^{1/2}e_I,\qquad V_I=\int_I|P_I|^2\,\dd t .
\end{equation}
We take $P_I$ as a Chebyshev interpolant, bound $e_I$ by an analytic envelope on
a Bernstein ellipse plus the node-value errors, and compute node values and $V_I$
in ball arithmetic \cite{Johansson2017}. The two one-sided fields are of order
one and their difference of order $10^{-12}$; ball arithmetic at 128 bits keeps
the cancellation inside the radii. The ingredients are listed in
\cref{tab:inputs}; all of them are computed for the same field, with all
rounding included.

\begin{table}[tbhp]
\caption{Inputs of the reliable bound and where their computation is described.}
\label{tab:inputs}
\centering\small
\begin{tabular}{lp{0.42\textwidth}l}
\toprule
Input & Content & Details\\
\midrule
$\overline\varepsilon_g$ & panel bounds \eqref{eq:panel}, self-interaction, local terms, tails \eqref{eq:tail} & \cref{sm:reliable-eps}\\
$\underline U$ & disk bounds \eqref{eq:disk} at certified point values & \cref{sm:reliable-mass}\\
$\overline C_g$ & collar checks and the supremum \eqref{eq:Cg} & \cref{sm:reliable-cg}\\
$\overline a_s,\underline b_s$ & inertia tests \eqref{eq:panel-tests} & \cref{sm:band-tests}\\
$R_{\mathrm{spec}}$, $I_k$ & scalar assembly, outward rounding & \cref{sm:assembly}\\
\bottomrule
\end{tabular}
\end{table}

\emph{Assembly and interpretation.} With $\underline D=\underline U-\sum_g
\overline C_g\overline\varepsilon_g>0$, \cref{thm:distance} gives
$\operatorname{dist}(\mu,\sigma(A(\beta)))\le R_{\mathrm{spec}}:=2\mu\sum_g
\overline C_g\overline\varepsilon_g/\underline D$; thus
$I_{\mathrm{spec}}=[\mu-R_{\mathrm{spec}},\mu+R_{\mathrm{spec}}]$ meets the
spectrum, and $I_k=[(\max\{0,\mu-R_{\mathrm{spec}}\})^{1/2},(\mu+R_{\mathrm{spec}})^{1/2}]$
contains the corresponding wavenumber. If in addition \eqref{eq:gap-condition}
holds, \cref{prop:discrete} shows that $I_{\mathrm{spec}}$ contains at least
one isolated eigenvalue of finite multiplicity, that is, a guided mode. We do not
claim that it is the only one, nor its index or multiplicity.

% Section 6. Budget 3.75 pages.
% Source: draft/bie-enclosure/s-results.tex (clean version); data from
% test/bie-bound/pipeline-v1 run-002, resolution-v1, example2-v1 run-015,
% gapedge-v1 run-006/007/011/012, test/fem-goerisch/pipeline-v1.
\section{Numerical results}\label{sec:results}
Each interval below comes from one run that starts from an empty output
directory and performs, in order, the gap classification, the determinant
search, the density reconstruction and field fit, the reliable evaluation of
all five groups, the mass bound, the trace constants, the assembly, and the
gap test, without reading earlier results. Geometry parameters that are not
binary64 numbers (such as $0.2$) are replaced by the nearest binary64 number,
and the intervals refer to this geometry. Timings are wall-clock on an Apple M3
Pro (11 cores, 18\,GiB), MATLAB R2023b and python-flint 0.9.0. The code, the
saved candidate data and an independent checker are in the release package
\todo{cite release}.

\begin{figure}[tbhp]
\centering
\begin{tikzpicture}
\begin{groupplot}[group style={group size=1 by 2, vertical sep=1.05cm},
  width=\textwidth, height=0.23\textwidth, axis equal image,
  ymin=-0.5, ymax=0.5, ytick={-0.5,0,0.5}, xtick={-3,...,3},
  ylabel={$y$}, unbounded coords=jump, tick label style={font=\scriptsize},
  label style={font=\scriptsize}, title style={font=\scriptsize}]
\nextgroupplot[xmin=-3.5, xmax=3.5, title={Example 1}]
  \foreach \w in {-2.5,-1.5,...,2.5} {
    \edef\temp{\noexpand\draw[gray, dashed] (axis cs:\w,-0.5) -- (axis cs:\w,0.5);}
    \temp}
  \pgfplotsinvokeforeach{-3,-2,-1}{\addplot[blue, thick, fill=blue!15] table[col sep=comma,
    x=x, y expr={\thisrow{cell}==#1 ? \thisrow{y} : nan}] {data/geometry.csv};}
  \addplot[red, thick, fill=red!15] table[col sep=comma,
    x=x, y expr={\thisrow{cell}==0 ? \thisrow{y} : nan}] {data/geometry.csv};
  \pgfplotsinvokeforeach{1,2,3}{\addplot[teal, thick, fill=teal!15] table[col sep=comma,
    x=x, y expr={\thisrow{cell}==#1 ? \thisrow{y} : nan}] {data/geometry.csv};}
\nextgroupplot[xmin=-3.5, xmax=3.5, xlabel={$x$}, title={Example 2 ($L_0=1$)}]
  \foreach \w in {-2.5,-1.5,-0.5,0.5,1.35,2.2,3.05} {
    \edef\temp{\noexpand\draw[gray, dashed] (axis cs:\w,-0.5) -- (axis cs:\w,0.5);}
    \temp}
  \pgfplotsinvokeforeach{-3,-2,-1}{\addplot[blue, thick, fill=blue!15] table[col sep=comma,
    x=x, y expr={\thisrow{cell}==#1 ? \thisrow{y} : nan}] {data/geometry2.csv};}
  \pgfplotsinvokeforeach{1,2,3}{\addplot[teal, thick, fill=teal!15] table[col sep=comma,
    x=x, y expr={\thisrow{cell}==#1 ? \thisrow{y} : nan}] {data/geometry2.csv};}
  % Fourth right rod: the third one shifted by one period L_+ = 0.85 (clipped).
  \addplot[teal, thick, fill=teal!15] table[col sep=comma,
    x expr={\thisrow{x}+0.85}, y expr={\thisrow{cell}==3 ? \thisrow{y} : nan}] {data/geometry2.csv};
\end{groupplot}
\end{tikzpicture}
\caption{Part of the strip in the two examples, on the same $x$ scale. Dashed
lines are the artificial walls; $y=\pm1/2$ is the seam. The left leads are the
same; in Example~2 the center cell is empty and the right lead has period
$L_+=0.85$ and a rotated rod. Lead rods have $\rho=17$, the star rod $\rho=4$.}
\label{fig:geometry}
\end{figure}

\subsection{Example 1: a star-shaped defect}\label{sec:ex1}
Let $d=1$, $\beta=0.5$, $L_0=L_\pm=1$, $\rho_{\mathrm e}=1$. The center cell
contains a star with $\rho=4$, and every lead cell a three-lobed rod with
$\rho=17$, centered in the cell,
\[
  \boldsymbol r_0(t)=0.2\begin{pmatrix}(1.5+0.3\cos5t)\cos t\\(1+0.3\cos5t)\sin t\end{pmatrix},
  \qquad
  \boldsymbol r_a(t)=\frac{0.2(1+a\cos3t)}{(1+a^2/2)^{1/2}}\begin{pmatrix}\cos t\\\sin t\end{pmatrix},
\]
with $a=0.12$ on the left and $a=0.20$ on the right (\cref{fig:geometry}). Thus $\rin=1$, $\rax=17$ and
$\lambda_0=1/68$.

\emph{Discretization.} The center curve has $512$ nodes, each lead curve $256$;
$M=48$ ($K=97$). The trial field uses $480$ proxies (\cref{sec:field}). The
reliable evaluation splits the center curve into $192$ panels, each lead curve
into $64$, each wall into $16$ and each seam into $8$, with $48$ Chebyshev nodes per
panel and $128$-bit balls. Both leads use $p=1$, with
$\|\mathsf P_\pm\|_F\le0.3815$, $0.3868$; $N_{\mathrm{cell}}$ is the smallest depth
at which the squared tail is at most $10^{-6}$ times the squared finite part
($52$--$58$). The mass bound uses $8$ disks inside the rods. The gap test uses
$169$ plane waves and $118$ phase panels on $[0,\pi]$; both lead rods are
symmetric under $y\mapsto-y$.

\emph{Candidate.} The floating-point gap classification gives the search window
$k\in[1.451,2.249]$. Chebyshev interpolation of $\det\mathsf A(k)$ on the whole
window stops at degree $32$ (last three coefficients below $7\times10^{-11}$
relative to the largest), with one near-real root in the Bernstein ellipse. Two
Newton steps give $\widehat k=1.558437894182088=438661269969907/2^{48}$; $\mu$ is
its exact square. The search takes $38$ matrix evaluations and $159$\,s.

\emph{Interval.} \Cref{tab:ex1} lists the five groups. The center rod
contributes $90\%$ of $Q$. There the collar is narrow ($h_g=0.969$) because the
concave points of the star have curvature about $64$; this gives the largest
$C_g$. The interpolation errors satisfy $e_I\le2.2\times10^{-17}$ on all panels, so
the reliable $\overline\varepsilon_g$ exceed the floating-point values by
only $0.04\%$. Increasing every $N_{\mathrm{cell}}$ by $10$ changes the group bounds by a
relative $4.7\times10^{-7}$ at most. The disks give $\underline U=1.23587$, and
\begin{align*}
  R_{\mathrm{spec}}&\le1.3690\times10^{-11},\\
  I_{\mathrm{spec}}&\subset[2.42872867000901,\,2.42872867003640],\\
  I_k&\subset[1.55843789417769,\,1.55843789418648],
\end{align*}
with outward rounding. All $472$ inertia tests pass, so neither lead has spectrum
for $k\in(1.318,2.375)$; the margins in \eqref{eq:gap-condition} are $0.6916$ and
$3.2118$. Hence $I_k$ contains at least one guided-mode wavenumber. The complete
run takes $827$\,s with a peak memory of $1.84$\,GiB; $528$\,s of this is the
reliable evaluation on $8$ processes (\cref{sm:results}).

\begin{table}[tbhp]
\caption{Example 1: trace constants and defect bounds of the five groups (upward
rounded). Lead and wall groups include all cells of both half guides.}
\label{tab:ex1}
\centering\small
\begin{tabular}{lccc}
\toprule
Group $g$ & $\overline C_g$ & $\overline\varepsilon_g$ & $\overline C_g\overline\varepsilon_g$\\
\midrule
center rod & $2.490254$ & $1.2616\times10^{-12}$ & $3.1417\times10^{-12}$\\
left rods & $1.181583$ & $5.0811\times10^{-14}$ & $6.0037\times10^{-14}$\\
right rods & $1.233941$ & $5.2316\times10^{-14}$ & $6.4555\times10^{-14}$\\
walls & $1.035553$ & $2.0346\times10^{-13}$ & $2.1070\times10^{-13}$\\
seams & $1.035553$ & $5.7837\times10^{-15}$ & $5.9894\times10^{-15}$\\
\midrule
$Q$ & & & $3.4830\times10^{-12}$\\
\bottomrule
\end{tabular}
\end{table}

\emph{Resolution.} We vary the field and the evaluation separately (complete
reruns, \cref{sm:results}). With $256$, $384$, $512$ center nodes (half as many per
lead curve) the full width $2R_{\mathrm{spec}}$ is $4.28\times10^{-6}$,
$9.87\times10^{-9}$, $2.74\times10^{-11}$, almost entirely through the center group;
$M=40$ or $56$ changes the width by about $1\%$. The five candidates differ from
each other by at most five ulps. For the same field, doubling all panels,
$64$ Chebyshev nodes or $192$-bit balls change each group bound by a relative
$3\times10^{-6}$ at most. The width therefore measures how well the reconstructed
field satisfies the interface conditions, not the accuracy of $\widehat k$ or of
the evaluation.

\emph{Comparison with a finite element enclosure.} The same problem was solved
independently with $P_3$ finite elements on a mesh with $h=1/48$ (about
$2.9\times10^4$ degrees of freedom per cell), a trial field built from a discrete
DtN map and tapered to zero over $13$ lead cells, and Goerisch's method
\todo{cite FEM route}. Under a spectral isolation condition, checked by separate
finite-dimensional tests, it gives $\mu\in[2.4287286691977,2.4287286736602]$ (width
$4.46\times10^{-9}$) together with uniqueness in this interval; truncation errors
are bounded reliably but rounding is only estimated. $I_{\mathrm{spec}}$ lies
inside. The two codes define the binary geometry slightly differently (by
$2^{-56}$ in one Fourier coefficient), so this is a numerical cross-check rather
than a statement about one operator. The FEM run takes $12.5$\,min plus
$54.3$\,min for the isolation condition.
\todo{FEM cost--width series over three mesh sizes (release R3).}

\subsection{Example 2 and guided modes near a gap edge}\label{sec:ex2}
The second example removes the center rod, uses different lead periods, and
breaks the reflection symmetry. The center cell $|x|<L_0/2$ is empty; the left
lead has $L_-=1$ and the rod $a=0.12$; the right lead has $L_+=0.85$ and the rod
$a=0.20$ rotated by $0.13$ (\cref{fig:geometry}). The right bands at $\pm\theta$
then differ by up to $3.7\times10^{-3}$, and the gap test covers the full range
$[-\pi/L_s,\pi/L_s]$ with $236$ panels per lead. The discretization is the same
as in Example~1. The center field is a finite Rayleigh sum (\cref{thm:center-rep}
with $\Omega_0=\emptyset$), which is exactly quasi-periodic, so the center rod and
center seam groups vanish identically, and the walls dominate $Q$ ($57\%$). The
floating-point common gap is about $k\in(1.441,2.396)$.

Increasing $L_0$ lowers the guided wavenumber toward the lower gap edge, and the
decay in the right lead slows down. \Cref{tab:ex2} reports complete reliable
runs for $L_0=1,1.6,1.9,2$; at each point the lowest confirmed root in the window
is used, not continued from the previous one. Along a floating-point continuation,
$\ln(1/\varrho(\mathsf P_+))\approx(1.94\text{--}2.04)\sqrt{\delta_{\mathrm e}}$, where $\varrho$ is the
spectral radius and $\delta_{\mathrm e}$ the distance to the band edge $1.441$, in line with the square-root behavior of
evanescent decay near a quadratic band edge. Correspondingly
$N_{\mathrm{cell}}\ln(1/\overline r_+)\approx51,49,48$, so the depth grows like
$\delta_{\mathrm e}^{-1/2}$. The width stays below $10^{-11}$
because the tail enters only through its bound. From $L_0=1$ to $2$ the right
depth grows fourfold but the total time only by $42\%$, since the search and the
gap test do not depend on $L_0$. Throughout,
$\|\mathsf P_\pm\|_2/\varrho(\mathsf P_\pm)\le1.021$ and $p=1$ suffices; the $p=2$
bound \eqref{eq:tail}, evaluated on the same states, differs by about $2\%$. For
$L_0=2$ the gap test window is $(1.45,2.37)$, the smallest inertia margin is
$0.197$, and the interval has margin $0.0161$ to the lower end; a second
candidate $\widehat k\approx2.2156$ appears in the window and is not treated.

\begin{table}[tbhp]
\caption{Example 2 for increasing center width $L_0$: interval $I_k$ (shared
leading digits, then outward-rounded last digits), distance $\delta_{\mathrm e}$ of
$\widehat k$ to the floating-point band edge, upper bound $\overline r_+$ of
$\|\mathsf P_+\|_2$, right depths $N_{\mathrm{cell}}$ (rods/walls/seams), $Q$, full
width of $I_{\mathrm{spec}}$ (bounds rounded upward), and time of one cold run of the release package. The depths are chosen so that the
squared tail is at most $\tau$ times the squared finite part, with
$\tau=10^{-6},10^{-2},0.1,0.1$; the tail always enters through its bound.}
\label{tab:ex2}
\centering\footnotesize\setlength{\tabcolsep}{3pt}
\begin{tabular}{cccccccc}
\toprule
$L_0$ & $I_k$ & $\delta_{\mathrm e}$ & $\overline r_+$ & $N_{\mathrm{cell}}$ & $Q$ & width & s\\
\midrule
$1$   & $1.89169614005[23,32]$ & $0.451$  & $0.3344$ & $52/51/52$ & $1.40\times10^{-13}$ & $3.09\times10^{-12}$ & $1076$\\
$1.6$ & $1.58144476391[49,70]$ & $0.140$  & $0.4924$ & $72/69/72$ & $4.13\times10^{-13}$ & $6.36\times10^{-12}$ & $1194$\\
$1.9$ & $1.48029280022[43,61]$ & $0.0393$ & $0.6825$ & $128/121/126$ & $3.87\times10^{-13}$ & $5.22\times10^{-12}$ & $1353$\\
$2$   & $1.45555716328[31,60]$ & $0.0146$ & $0.7953$ & $208/196/206$ & $6.24\times10^{-13}$ & $8.16\times10^{-12}$ & $1532$\\
\bottomrule
\end{tabular}

\end{table}

% Section 7. Budget 0.25 pages.
\section{Conclusions}\label{sec:conclusion}
A field that solves the Helmholtz equation exactly in each homogeneous piece of
the strip gives a bound on the distance to the spectrum in terms of its interface
jumps, with constants that depend only on the geometry. With a boundary integral
candidate, this yields rounding-inclusive enclosures of guided-mode eigenvalues
of width $10^{-12}$--$10^{-11}$ in $14$--$25$ minutes, including the gap test,
also when the mode approaches a gap edge. The width is governed by how well the
reconstructed field matches across the rod boundaries, not by the evaluation.
Uniqueness and the index of the enclosed eigenvalue are not addressed; in the
finite element comparison they come from a separate isolation condition.

% Acknowledgments with the disclosure of AI use required by the SIAM
% Publications Editorial Policy on Artificial Intelligence, Version 2.0
% (effective May 2026): usage reported in the Acknowledgements with the
% statement "The authors assume responsibility for all content." The items
% follow the policy's areas for diligence. The policy does not ask for model
% versions, so only the tool families are named.
\section*{Acknowledgments}
\emph{Use of AI tools.} AI tools (Claude by Anthropic and GPT models by OpenAI)
were used substantially in this work between May and November 2026.
\begin{itemize}
\item \emph{Theorems and proofs.} GPT models drafted and audited proofs of the
  results in \cref{sec:model,sec:reduction,sec:discretization,sec:bound} and in
  the supplementary materials, including \cref{thm:complement},
  \cref{prop:trace}, and \cref{thm:distance}, as well as the error bounds of the
  rounding-inclusive evaluation.
\item \emph{Code and experiments.} GPT models wrote the MATLAB code for the
  discretization and the root search, the Python code for the ball-arithmetic
  evaluation and the gap test, and the experiment drivers, and ran the
  experiments. Claude checked the outputs against the experiment specifications
  and reran the verification layer of the release package. All code and the
  data that define the reported fields are public in that package
  \todo{cite release}.
\item \emph{Figures and tables.} The figure and all tables are produced by
  scripts from computed data; no figure was generated directly by an AI tool.
\item \emph{Literature and text.} The AI tools were used to search for related
  literature and to draft the related-work paragraphs. Claude drafted and
  revised the text of this paper and the supplementary materials; GPT produced
  referee-style reports on earlier drafts, which guided revisions.
\item \emph{Checking by the authors.} \todo{State what the authors checked
  personally: proofs checked line by line, code tested, every cited work read
  and its summary checked, numbers reproduced from the release package.}
\end{itemize}
The authors assume responsibility for all content.


\bibliographystyle{siamplain-nonote}
\bibliography{../../ref/reference}
\end{document}
